\labsection{2}{File scanning and directory utilities}{sec:lab02}
\begin{labproject}{Build small Unix-style text and directory tools}
\textbf{Coverage.} Chapter 2.\\
\textbf{Source files.} \texttt{source\_code/Lab02/task\_1.c}--\texttt{task\_3.c}.\\
\textbf{Goal.} Practice robust read loops, pathname arguments, directory traversal, and line-oriented text searching.

\medskip\noindent\textbf{Task A --- Count spaces safely.}\par
Replace a \texttt{while (!feof(fp))} loop with a read-result loop:
\begin{lstlisting}[style=osc]
int ch, spaces = 0;
while ((ch = fgetc(fp)) != EOF) {
    if (ch == ' ') spaces++;
}
\end{lstlisting}

\medskip\noindent\textbf{Task B --- Fix the directory listing.}\par
The supplied \texttt{task\_2.c} closes the directory inside the traversal loop. Move \texttt{closedir(dp)} after the loop and add an \texttt{opendir()} failure check.

\medskip\noindent\textbf{Task C --- Mini grep.}\par
The supplied program reads each line with \texttt{fgets()}, removes the newline if present, uses \texttt{strstr()} to search for a substring, and prints matching lines with line numbers.

\begin{debugtask}
Compile every Lab 2 source with warnings enabled, for example \texttt{gcc -Wall -Wextra file.c}. Record each warning/error and fix the program rather than suppressing diagnostics.
\end{debugtask}
\end{labproject}

\begin{labdeliverable}
Submit corrected source for all three tasks and a one-page note explaining why successful-read loops and acquire/use/release resource discipline matter in OS programming.
\end{labdeliverable}
